\begin{figure}[H]
\centering
\begin{tikzpicture}[ancho/.style={minimum width=3cm}]
    \node[actor, ancho] (saved) at (0,0) {Se guarda\\un cambio};
    \node[actor, ancho] (manual) at (4.2,0) {Publicar\\a mano};
    \node[datos, minimum width=2.8cm] (logdb) at (8.6,0) {SQLite\\programación\\y estado};

    \node[servicio, ancho] (autosnap) at (4.2,-2.9) {Respaldo\\automático};
    \node[servicio, ancho] (copyfiles) at (8.6,-2.9) {Copia de JSON,\\Markdown e imágenes};
    \node[servicio, ancho] (build) at (13,-2.9) {Construcción\\del sitio};
    \node[interfaz, ancho] (published) at (13,-5.3) {Sitio\\publicado};

    \draw[flechaauto] (saved) -- node[etq, pos=0.4, left=6pt] {a los pocos\\segundos} (autosnap);
    \draw[flechaadm] (manual) -- (autosnap);
    \draw[flechaauto] (logdb) -- node[etq, pos=0.5, right=10pt] {llega la\\fecha} (autosnap);
    \draw[flecha] (autosnap) -- (copyfiles);
    \draw[flecha] (copyfiles) -- (build);
    \draw[flecha] (build) -- (published);
    \draw[flecha] (build.north) |- node[etq, pos=0.75, above=2pt] {registra\\el estado} (logdb.east);
\end{tikzpicture}
\caption[Proceso de publicación del sitio informativo]{Proceso de revisión, respaldo, construcción y publicación del sitio informativo}
\label{fig:cms_proceso_publicacion}
\end{figure}
